\section{Conclusion}
\label{sec:conclusion}

The geometry of the polarization ellipse 
that defines the Stokes parameters is elegantly connected
to that of the coherency matrix via the Pauli matrices.
%
Geometry also relates linear transformations of the electric field 
to their impact on the Stokes parameters, leading to a physically meaningful
and broadly applicable classification of polarimetric transformations.
%
%Following a polar decomposition into Hermitian and unitary matrices, the axis-angle representation of these transformations promotes an intuitive geometric interpretation of their effect on the Stokes parameters.
%
% Well-studied Lorentz boosts and Euclidean rotations of the Stokes parameters directly map to physical properties of a system known as diattenuation and retardance, respectively.
%
The properties of these transformations simplify the analysis of the elements of the signal path and yield both theoretical insights and practical guidelines that can be applied during calibration.

These guidelines include criteria for evaluating and selecting a model of the instrumental response that performs well during numerical analysis of experimental data.
%
%
Among those considered,  the model introduced by \citet{bri00} satisfies all of the selection criteria for single-antenna observations of a point source.
%
In this physically-motivated and self-consistent model, there is a one-to-one mapping between parameters and linear transformations, which are described
using equations that remain numerically stable during modeling.

This introduction focuses primarily on the mathematical and conceptual foundations of polarimetry, with emphasis on the numerical analysis of radio astronomical observations.  For a broader review of both the physical processes that produce polarized radiation and the astrophysical phenomena that can be studied through polarimetry, the review by \cite{tri14} is exceptionally thorough and insightful.

Radio polarimetry is a vibrant field of research that continues to yield transformative insights and discoveries in astrophysics, driven by both technological advances and innovative techniques.  For example, the relatively recent discovery of polarization closure traces \citep{bp20,snt22} enabled the first images of magnetic field structure near the M87 black hole using very long baseline interferometry \citep{2021ApJ...910L..13E}.
%
Novel techniques of high-fidelity polarimetry led to the discovery of a surprisingly high degree of linear polarization in low-frequency radio images of solar bursts \citep{dkom25}.
%
At even higher time resolution, measurements of the polarized radiation from pulsars, magnetars and fast radio bursts continue to reveal new insights into their nature and extreme environments \citep[e.g.,][]{crh+06,2013Natur.501..391E,2018Natur.553..182M,ptv+22}.

Current and future scientific goals, such as detecting the neutral Hydrogen emission from the epoch of reionization \citep[e.g.,][]{2018ApJ...867...15T,2020MNRAS.493.1662M} and understanding the origin and evolution of interstellar and intergalactic magnetic fields \citep[e.g.,][]{2017ARA&A..55..111H}, motivate the development of next-generation facilities and methods of polarimetric analysis, such as Faraday tomography \citep[e.g.,][]{2000JGR...10516063G,2017A&A...597A..98V}.
%
Furthermore, the unprecedented sensitivity, resolution, and fields of view of new observatories necessitate further research and development of novel methods of direction-dependent interferometric calibration \citep[e.g.,][]{2015MNRAS.449.2668S} and correction of ionospheric Faraday rotation \citep[e.g.,][]{2019A&A...622A...5D}.
%
Improved polarimetric fidelity also has the potential to increase the sensitivity of pulsar timing array experiments to the low-frequency stochastic gravitational wave background \citep[e.g.,][]{gcv+23,rvg+24,gvc+25}.

Given the breadth and potential impact of scientific knowledge that remains to be discovered through polarimetry, I hope that this introduction might help others to further pursue the subject.
